% Part: computability
% Chapter: computability-theory
% Section: reducibility

\documentclass[../../../include/open-logic-section]{subfiles}

\begin{document}

% SEGMENT OLP-0243-S01

\olfileid{cmp}{thy}{red}
\olsection{குறைத்தன்மை}

\begin{explain}
கணிக்கத்தக்கவாறு எண்ணிடத்தக்கதும் கணிக்கத்தகாததுமான ஒரு கணமாவது,
$K_0$, உள்ளது என்று இப்போது அறிவோம். பிற கணங்களும் உள்ளன என்பது
தெளிவாக இருக்க வேண்டும். ஒவ்வொரு முறையும் முதல் கொள்கைகளுக்குத்
திரும்ப வேண்டியிராமல், மற்ற கணங்களுக்கு இப்பண்புகள் உள்ளன என்று
காட்டுவதற்குக் குறைத்தன்மை முறை வலிமையான வழியை வழங்குகிறது.

பொதுவாக, கணம் $A$ ஐ கணம்~$B$ க்கு ``குறைத்தல்'' என்பது,
!!{element}s கணம்~$B$ இல் உள்ளனவா இல்லையா என்பதற்கான விடைகளை,
!!{element}s கணம்~$A$ இல் உள்ளனவா இல்லையா என்பதற்கான விடைகளாக
மாற்றும் ஒரு முறையாகும். ``பலவற்றுக்கு ஒன்று குறைத்தன்மை''
என்ற கருத்தில் நாம் கவனம் செலுத்துவோம்; ஆனால் வெவ்வேறு பண்புகளைக்
கொண்ட பல குறைத்தன்மைக் கருத்துகளும் உள்ளன. கணிப்புச் சிக்கற்பாட்டை
ஆராய்வதிலும் குறைத்தன்மைக் கருத்துகள் மையமானவை; அங்கு செயற்றிறன்
சிக்கல்களையும் கருத்தில் கொள்ள வேண்டும். எடுத்துக்காட்டாக, ஒரு கணம்
NP இல் இருந்து, ஒவ்வொரு NP சிக்கலையும் அதற்குக் குறைக்க முடிந்தால்,
அது ``NP-முழுமையானது'' எனப்படும். கீழே விவரிக்கப்படும்
குறைத்தன்மையைப் போன்ற ஒரு கருத்தே அங்கு பயன்படுகிறது; கூடுதலாக,
அக்குறைத்தலைப் பல்லுறுப்புக்கோவை நேரத்தில் கணிக்க முடிய வேண்டும்.

குறைத்தல் என்ற கருத்தை நாம் ஏற்கெனவே மறைமுகமாகப் பயன்படுத்தியுள்ளோம்.
கணம்~$K$ ஐ
\[
K = \Setabs{x}{\cfind{x}(x) \fdefined},
\]
அதாவது $K = \Setabs{x}{x \in W_x}$ என வரையறுக்க. நிற்றல் சிக்கல்
தீர்க்கவியலாதது என்பதற்கான நமது நிரூபணம்
(\olref[hlt]{thm:halting-problem}), $K$ கணிக்கத்தகாதது என்பதை
மிக நேரடியாகக் காட்டுகிறது. $K_0$ என்பது
\[
K_0 = \Setabs{\tuple{e, x}}{\cfind{e}(x) \fdefined },
\]
அதாவது $K_0 = \Setabs{\tuple{x,e}}{x \in W_e}$ என்ற கணம் என்பதை
நினைவுகூர்க. $K$ கணிக்கத்தகாதது என்பதற்கான எந்த நிரூபணத்தையும்
$K_0$ கணிக்கத்தகாதது என்பதற்கான நிரூபணமாக எளிதில் நீட்டிக்கலாம்:
$K_0$ கணிக்கத்தக்கதாக இருந்தால், !!a{element}~$x$ கணம் $K$ இல்
உள்ளதா இல்லையா என்பதை, சோடி $\tuple{x, x}$ ஆனது $K_0$ இல்
உள்ளதா இல்லையா என்று கேட்பதன்மூலம் தீர்மானிக்கலாம். $x$ ஐ
$\tuple{x, x}$ க்கு அனுப்பும் சார்பு~$f$, $K$ இலிருந்து $K_0$
க்கான ஒரு \emph{குறைத்தலுக்கு} எடுத்துக்காட்டு.
\end{explain}

\begin{defn}
$A$ உம் $B$ உம் இயல் எண்களின் கணங்கள் ஆகட்டும். ஒவ்வோர் இயல்
எண்~$x$ க்கும்
\[
x \in A \quad \text{எனில், எனில் மட்டுமே} \quad f(x) \in B.
\]
எனும் நிபந்தனை நிறைவேறினால், கணிக்கத்தக்க சார்பு
~$f\colon \Nat \to \Nat$ ஆனது $A$ இலிருந்து~$B$ க்கான ஒரு
\emph{பலவற்றுக்கு ஒன்று குறைத்தல்}. இத்தகைய குறைத்தல்~$f$ இருந்தால்,
$A$ ஆனது~$B$ க்கு \emph{பலவற்றுக்கு ஒன்று குறைக்கத்தக்கது} என்கிறோம்;
இதை $A \leq_m B$ என்று எழுதுகிறோம். $A$ ஆனது $B$ க்கும், $B$ ஆனது
$A$ க்கும் பலவற்றுக்கு ஒன்று குறைக்கத்தக்கவையாக இருந்தால், அவை
\emph{பலவற்றுக்கு ஒன்று சமானமானவை} எனப்படும்; இதை
$A \equiv_m B$ என்று எழுதுகிறோம்.
\end{defn}

மேலுள்ள வரையறையிலுள்ள சார்பு $f$ உட்செலுத்தலாகவும் இருந்தால்,
$A$ ஆனது~$B$ க்கு \emph{ஒன்றுக்கு ஒன்று குறைக்கத்தக்கது} எனப்படும்.
கீழே விவரிக்கப்படும் பெரும்பாலான குறைத்தல்கள் இந்த வலுவான
நிபந்தனையையும் நிறைவேற்றுகின்றன; ஆனால் நாம் இவ்வுண்மையைப்
பயன்படுத்தமாட்டோம்.

\begin{digress}
ஒன்றுக்கு ஒன்று குறைத்தன்மை உண்மையிலேயே பலவற்றுக்கு ஒன்று
குறைத்தன்மையைவிட வலுவான நிபந்தனை; ஆனால் இது எவ்விதத்திலும்
வெளிப்படையானதல்ல. வேறுவிதமாகக் கூறினால், $A$ ஆனது~$B$ க்கு
பலவற்றுக்கு ஒன்று குறைக்கத்தக்கதாகவும், ~$B$ க்கு ஒன்றுக்கு ஒன்று
குறைக்கத்தகாததாகவும் இருக்கும் முடிவிலாக் கணங்கள் $A$,~$B$ உள்ளன.
\end{digress}

\end{document}
